\documentclass[11pt]{article}
\usepackage[letterpaper,margin=0.7in]{geometry}
\usepackage[T1]{fontenc}
\usepackage{lmodern}
\usepackage{amsmath,amssymb,tikz,array,fancyhdr}
\renewcommand\familydefault{\sfdefault}
\setlength\parindent{0pt}
\setlength\parskip{5pt}
\setlength{\headheight}{15pt}
\pagestyle{fancy}\fancyhf{}
\lhead{Week 3 / Shuffle machines / Adult return-visit guide}
\fancyfoot[L]{Bellingham Math Circle / Week 3 / F03-RV-FAC-v1}
\fancyfoot[R]{\thepage}
\renewcommand\headrulewidth{0.4pt}
\newcommand\lead[1]{\textbf{#1}\par}
\begin{document}
\lead{Mathematical overview: three ways to constrain a shuffle}
For distinct ordered cards, an inversion is a pair whose larger number appears first. An adjacent swap changes the inversion count by exactly one. Repeatedly swapping adjacent inverted pairs reaches the home order, so the exact minimum sorting distance is the inversion count. For $n$ cards the reverse order requires $n(n-1)/2$ swaps; nonadjacent exchanges are a different rule.

If the only moves are rotating the last card to the front and reversing the entire row, the reachable orders are precisely the $n$ cyclic readings of the starting row and the $n$ cyclic readings of its reverse, for $n\ge 3$ distinct labels. Thus four cards have eight reachable orders out of 24, five have ten out of 120. The circular unordered-neighbor relation is preserved. Repeated labels alter the counts.

A machine run twice is the square of its permutation. A target has an identical-machine two-turn reconstruction exactly when the number of its cycles of each even length is even. Odd cycles have roots individually; equal-length even cycles can be paired and interlaced. One swap alone has no root, two disjoint swaps have a four-cycle root, and a three-cycle has a root. This is an adult theorem map; children can build and run small candidate machines without cycle notation.

Experiments find candidate routes and machines. Inversion counting explains a minimum, circular neighbors explain nonreachability, and cycle splitting explains square-root obstructions. A handful of unsuccessful machine drawings is not a proof of impossibility.

\lead{Entries, materials and a first route}
Shared Grades 2--5: sorting and physical rotation/reversal can suit younger children with adult reading and a short row. Reconstructing a machine from two turns suits Grades 4--5, or younger children already fluent in the base arrow-machine action. Counting/order through six is enough for the concrete core. Holding two represented machine states at once is a readiness gate; use objects and an unchanged target instead of mental bookkeeping.

Per pair: six distinct numbered slips or labeled blocks, pencil, a visible home row and a blank arrow machine. Slips should be about 20--25 mm wide so six fit across a Letter page; ordinary full-size playing cards can instead be laid out freely on the table. Page 1 supplies six 25 mm cutout cards: print a second copy of Page 1 per pair and cut them before the visit, leaving the recording page intact. The five-slot working mat has 25 mm squares; other machine and order diagrams are compact records, and loose cards operate on the table. For five pairs prepare 30 slips plus a demonstration set. At the upper trio rotate builder, checker and recorder. Preserve fixed three adult-led tables.

Choose one problem for 20--40 minutes, with free handling and a 2--4 minute whole-group launch. Keep unused pages for later visits. The lower table may stay with sorting or physical move trials; no need to force square roots onto every child. Adults may read or hold the comparison state; children should choose the swaps or candidate arrows.

\lead{Launch together}
Show only the chosen action on an unrelated short row: exchange two neighbors, move the last card to the front, reverse a whole row, or run the same arrow machine twice. Let children handle the slips first. For two-turn reconstruction retain the target row unchanged and use a separate working row. Show input, first-turn intermediate, and second-turn output before asking for a machine; do not demonstrate one of the target roots.
\newpage
\lead{Problem 1: shortest adjacent sorting}
\textbf{Entry and timing.} Pairs exchange neighboring cards physically and retain a move word, identifying a swap by the two positions. Allow 20--35 minutes to improve routes and invent hard rows. A younger entry uses three or four labels, with the home order visible.

\textbf{Exact method and proof.} Count each out-of-order pair once. In $3,1,4,2$ the inverted pairs are $(3,1)$, $(3,2)$, $(4,2)$, so the minimum is three. One route is $3142\to\allowbreak 1342\to\allowbreak 1324\to\allowbreak 1234$. An adjacent exchange changes the order of only that neighboring pair, leaving its comparisons with all other cards unchanged. Therefore each swap changes the inversion count by one; at least that many swaps are needed. If a row is not sorted, some neighboring pair is inverted; swapping such a pair decreases the count by one. This reaches zero in exactly the lower-bound number of swaps.

\textbf{Printed rows.} The rows 21354, 31425 and 53421 require respectively 2, 3 and 9 adjacent swaps. The hardest five-card row is 54321, requiring 10. The three-card action example exchanges positions 1--2 of 231 to make 321; it demonstrates a legal move, without promising that move is helpful for sorting.

\textbf{Questions and hints.} ``Which two cards changed their relative order?'' If children count only adjacent reversed pairs, ask them to trace one large card past all smaller cards to its right. Do not require them to count inversions concurrently with moving; the preserved move word can be counted afterward.

\textbf{Return visits.} Make the hardest possible row and explain its bound: all $\binom n 2$ pairs are inverted in a reverse row. Seek a row with each possible distance from zero to that maximum. Change to swapping any two positions only as a distinct later rule; the inversion distance then ceases to be the minimum.

\textbf{Limits and novelty.} Distinct labels, fixed increasing home order, one neighboring exchange per move. Sorting distance is new; base machine periods, inverse machines and perfect-shuffle orders do not answer it.
\newpage
\lead{Problem 2: rotate or reverse}
\textbf{Entry and timing.} Children operate the two whole-row moves with slips, allowing 20--35 minutes. Keep one starting row unchanged as the reference. A younger entry is trying to make a partner's target; the older route is determining all reachable orders and diagnosing an impossible target.

\textbf{Exact result.} From $1234$ the eight reachable orders are $1234,2341,3412,4123$ and $4321,3214,2143,1432$. The other sixteen orders of those four distinct cards cannot be made. From $12345$ the ten are the five rotations of $12345$ and of $54321$. A rotation reads the same circular neighbor order from another start; reversal reads it in the opposite direction. Every allowed move keeps every unordered circular neighbor pair, including last--first. Every cyclic reading in either direction is attainable, so these are exactly all reachable rows.

\textbf{Printed targets.} From 1234, 3412 is reachable by two last-to-front moves; 2143 is reachable by two last-to-front moves followed by reversal. The targets 1324 and 2413 are impossible: their circular neighbor sets differ from the starting set. The eight reachable orders listed above fill the collection; ten printed blank records allow children to keep and then reject duplicates.

\textbf{A useful bridge.} After children have encountered an impossible target, place the row around a circle, joining last to first. Let them compare which labels touch on the starting circle and the target. This recording representation answers a question they already have; do not begin by printing the invariant as a recipe.

\textbf{Questions and hints.} ``Which two labels are still neighbors when the row wraps around?'' If they reverse only two cards, demonstrate reversing the entire row on an unrelated input. Let the partner check each legal move. Cards crossing during a physical reversal are not additional swaps; a reversal is one whole-row move under this problem's rule.

\textbf{Continuations.} Which unreachable target can be reached after allowing one adjacent swap in addition? With duplicate labels, some rotations coincide; investigate how the reachable count changes without applying $2 n$ blindly. Seek shortest rotate/reverse move words for reachable targets; reachability and optimal distance are different questions.

\textbf{Limits and novelty.} Exactly the two allowed move types; distinct labels for the counts above. The circular adjacency is unordered, because reversing is legal. This differs from unrestricted composition/inverse machine questions in the base packet.
\newpage
\lead{Problem 3: one machine, two turns}
\textbf{Entry and timing.} Use this with children who can retain the arrow-machine rule after a normal launch. Every top slot has one arrow out and every bottom slot one arrow in. Run all blocks simultaneously down their arrows, then move the whole bottom row straight back up. A second turn uses exactly the same arrows. Allow 25--45 minutes with a later visit for explanation.

\textbf{Small constructions.} A target swapping 1 and 2 while fixing 3 has no root. A target $1\to 2\to 3\to 1$ can be made by twice running $1\to 3\to 2\to 1$. A target swapping 1 with 2 and 3 with 4 is made by twice running $1\to 3\to 2\to 4\to 1$. Directly check every labeled block against the unchanged target; the intermediate row makes the verification recoverable.

\textbf{Printed two-pass targets.} These are output card rows, not input-to-output arrow lists. For output 231, use arrows $1\to2,2\to3,3\to1$: the first output is 312 and the second is 231. Output 213 is impossible. For 2143 use arrows $1\to3,2\to4,3\to2,4\to1$: first output 4312, second 2143. For output 231564 use arrows $1\to2,2\to3,3\to1$ and $4\to5,5\to6,6\to4$: first output 312645, second 231564. Other roots may work. The two-slot convention example swaps cards 7 and 8, then the same swap restores 7,8.

\textbf{Why the criterion works.} Squaring an odd cycle follows every second card and remains one cycle of the same odd length; stepping $(k+1)/2$ around a target odd $k$-cycle gives a root because twice that step is $1$ modulo $k$. Squaring a cycle of even length $2 k$ separates its alternating positions into two $k$-cycles. Hence target cycles of an even length must occur in pairs. Conversely, interlace two same-length cycles to make a cycle twice as long. Odd cycles can be handled separately, including fixed points. These constructions prove both necessity and sufficiency, assuming permutations of the same labeled slots.

\textbf{Hints.} ``Where does this one block go on turn one, and then on turn two?'' If the child changes arrows between turns, retain the drawn machine and ask the partner to enforce it. For two disjoint swaps, offer a four-cycle only after substantial attempts; for a single swap, invite them to list the possible loop shapes of three or four slots. No algebraic notation is needed.

\textbf{Return visits.} Pair two target three-cycles into one six-cycle root, or build roots for a five-cycle. Ask whether a target can have several roots: identity has many self-inverse roots, so reconstruction need not be unique. Enumeration checks: among all targets, roots exist for 3 of 6 three-slot permutations, 12 of 24 four-slot permutations, and 270 of 720 six-slot permutations. These are adult finite-check facts, not an elementary requirement to list hundreds of machines.

\lead{Evidence and status}
All three tasks change the question beyond base machine periods, inverses, composition and perfect shuffles. Exact distances/reachable sets/roots are independently checked in supplied code and fresh math review. \emph{Math Circle by the Bay}, Preface viii--x (PDF 9--11), supports deeper revisits and adaptable pacing; the particular tasks are new designs. Student and adult pages are \textbf{unpiloted}; physical procedure rehearsal is \textbf{untested}. Record concrete examples used and whether the children retained the move rule without the adult choosing their next attempt.

\end{document}
